\begin{table}[H]
\small
\setlength{\tabcolsep}{3pt}

\caption{\label{tab:rent-mapping-robustness}Robustness to rent treatment: inflation gap (in points, 2022)}
\centering
\begin{tabular}[t]{lrrr}
\toprule
Country & All rents & Effective rents & Excluding rents\\
\midrule
Austria & -0.2 & -0.7 & 0.1\\
Belgium & 2.0 & 0.5 & 2.4\\
Croatia & 0.3 & 0.3 & 0.4\\
Cyprus & 0.6 & 0.5 & 0.7\\
Estonia & 4.5 & 4.6 & 4.7\\
\addlinespace
Finland & -1.1 & -2.4 & -0.9\\
France & 0.1 & -0.7 & 0.2\\
Germany & 0.9 & -0.6 & 1.4\\
Greece & 1.1 & 0.5 & 1.3\\
Ireland & 2.8 & 2.9 & 3.0\\
\addlinespace
Italy & 1.5 & 1.5 & 1.6\\
Latvia & 4.6 & 4.6 & 4.8\\
Lithuania & 3.3 & 3.2 & 3.4\\
Luxembourg & 0.2 & -0.8 & 0.2\\
Malta & 0.2 & 0.2 & 0.3\\
\addlinespace
Netherlands & 1.8 & -0.4 & 2.7\\
Portugal & 0.2 & -0.5 & 0.3\\
Slovakia & 1.1 & 0.7 & 1.2\\
Slovenia & 0.5 & 0.7 & 0.5\\
Spain & 0.8 & 0.3 & 1.0\\
\addlinespace
Euro area & 0.9 & 0.0 & 1.2\\
\bottomrule
\end{tabular}

\begin{flushleft}
\footnotesize{\emph{Notes:} All scenarios use national-quintile Eurostat HBS baskets, relative-expenditure weights and chained Laspeyres indices. The first pools HBS 041+042 and maps it to HICP 041; the second uses HBS 041 only; the third excludes HBS 041 and 042 and HICP 041 and renormalises. Gap denotes first-quintile minus fifth-quintile inflation, in percentage points. Euro-area indices aggregate national indices with annual HICP country weights. Italy uses the custom national HBS input; Slovakia uses the 2015 HBS wave.}
\end{flushleft}
\end{table}
